\documentclass{article}
\usepackage[margin=2cm]{geometry}
\usepackage{amsmath}
\usepackage{graphicx}
\usepackage{booktabs}
\usepackage[utf8]{inputenc}
\usepackage{ragged2e}
\setlength{\emergencystretch}{3em}
\begin{document}
\sloppy

R-symmetry. As argued in, the field  is the top component of the global current
supermultiplet. Therefore a deformation by Ogo will break superconformal symmetry but
preserve all Poincare supersymmetry . However, deformation by Ogo alone is inconsistent.
Poincare supersymmetry demands that we also turn on the scalar masses O. Moreover,
this operator to be Og3.


\subsection*{1 Euclidean theory and BPS solutions}


In this section we will obtain the six-dimensional holographic dual of a mass deformation
of a 5D SCFT on S5. Such a dual is given by S5-sliced domain wall solutions of mattercontinue the Lorentzian signature gauged supergravity outlined above to Euclidean signature,
which has subtleties for both the scalar and fermionic sectors. Once the Euclidean theory
is obtained, we turn on relevant scalars necessary to support the domain wall. As discussed
in the previous section, at least three scalars must be turned on to obtain supersymmetric
solutions. The ansatz for the domain wall solutions takes the following form


\begin{equation}
d s ^ { 2 } = d u ^ { 2 } + e ^ { 2 f ( u ) } d s _ { S ^ { k } } ^ { 2 }, ~ ~ ~ ~ ~ ~ ~ \sigma = \sigma ( u ), ~ ~ ~ ~ ~ ~ ~ ~ \phi ^ { i } = \phi ^ { i } \{ u \}, ~ ~ i = 0, 3 ~
\end{equation}


with the remaining fields set to zero. Next we will obtain a consistent set of BPS equations
on the above ansatz, and then solve them numerically. When solving them, we will demand
as an initial condition that for some finite u the metric factor e2f vanishes, so that the.
geometry closes off smoothly.


\subsection*{1.1Euclidean action}


The Euclidean action may be obtained from the Lorentzian one by first performing a simple
Wick rotation of Lorentzian time t -> -ix6. This makes the spacetime metric negative definite, since the metric in the Lorentzian theory was taken to be of mostly negative signature.
However, we will choose to work with the Euclidean theory with positive definite metric.
Making this modification involves a change in the sign of the Ricci scalar. Then noting that
the Euclidean action is related to the Lorentzian action by exp (iSLor) = exp (-sEuc), the
final result of the Wick rotation is the following Euclidean action,


\begin{equation}
S _ { 8 D } = \frac { 1 } { 4 \pi G _ { 6 } } \int d ^ { 5 } x ~ \sqrt { G } \mathcal { L } ~, ~ ~ ~ ~ \mathcal { L } = \left ( - \frac { 1 } { 4 } R + \partial _ { \mu } \sigma \partial ^ { \mu } \sigma + \frac { 1 } { 4 } G _ { i j } ( \phi ) \partial _ { \mu } \psi ^ { i } \partial ^ { \mu } \phi ^ { j } + V \{ \sigma, \phi ^ { i } \} \right ) \quad ( 2 )
\end{equation}


where the spacetime metric G is positive definite and G6 is the six-dimensional Newton's
constant. By abuse of notation, Gij(\$) with indices refers to the metric on the scalar manifold

\end{document}
